\chapter{通货膨胀：每天从你口袋里拿走一点的那个人}
\chaptermeta{04}

\begin{ledetext}
上一章拆利息的时候，"通胀预期"那一块被放过去了。这一章回来算账。
\end{ledetext}

\section{赵敏芝的账本}

赵敏芝 58 岁，沈阳铁西人，退休前在一家轴承厂做了二十六年会计。退休金每月 3870，平时帮女儿带外孙，每周三上午去兴顺街的菜市场。

买菜的小票她还留着，按月夹在一个塑料夹子里。

所以她的账记得比统计局细：上个月排骨 21 块 8 一斤，这个月 24 块 5。

她的结论是，什么都在涨。

同一时期，全国 CPI 同比印出来是零点几。

两边都没撒谎。她们说的不是同一件事。

先把词抠清楚。\term{通货膨胀}{inflation}\index{061@通货膨胀}不是"某样东西涨价"，是\textbf{一般物价水平的持续、普遍上涨}。三个词都得较真：一般是整体不是某一样，持续是一直在涨不是涨完就停，上涨是水平在爬不是水平很高。

排骨那 2 块 7，动的是\term{相对价格}{relative price}\index{069@相对价格}。

相对价格是整个经济体的喊话系统。猪肉贵了，是市场扯着嗓子喊"多养点猪"；手机降价了，是市场在说"这条产线卷完了，人去干别的"。这些喊话指挥着几亿人往哪儿投钱、往哪儿投时间。

通胀不喊话。通胀干的是另一件事：把量身高的那把尺子改短。

\begin{pullquote}{}
相对价格在传递信息，一般物价水平在改变尺子。混成一件事，你会一边骂排骨，一边完全看不见真正在动的东西。
\end{pullquote}

\section{一个不存在的家庭}

CPI 的算法朴素得有点土。统计部门先假想一个"平均家庭"，把这家人一个月买的东西列成清单：粮油肉蛋、房租水电、公交地铁、理发看病、孩子的补习班，每项定个权重。然后每月跑几十万个采价点抄价，加权算总账单变了多少。

麻烦在于，这个"平均家庭"现实里一户都不存在。

\begin{egbox}{先算一笔账}
赵敏芝和老伴一个月开销 6480 块：食品烟酒 2130，居住 1180，医疗保健 940，教育文化娱乐 610，交通通信 520，其他用品及服务 460，生活用品 380，衣着 260。

这个月食品涨 5.8\%，医疗涨 2.4\%，交通通信因为油价回落降 4.1\%，其余没动。

2130×1.058 = 2253.5（多花 123.5），940×1.024 = 962.6（多花 22.6），520×0.959 = 498.7（少花 21.3）。总账单从 6480 变成 \textbf{6604.8，涨 1.93\%}。

而她走出菜市场时脑子里只有一个数：\textbf{5.8\%}。她记住的是权重最大的那次刺痛，统计算的是整张账单。

\end{egbox}

官方数字和体感对不上，有四条原因，每条都成立。

\begin{flowlist}
  \flowitem{01}{每个人的篮子不一样}{赵敏芝每月 940 块花在医药和门诊上，占她开销的 14.5\%。全国口径里医疗保健的权重要低得多。她的通胀率是她那个篮子的通胀率，跟"平均篮子"没有理由相同。她女儿家又是另一副样子：两个孩子，一个上小学一个上幼儿园，课外班加儿科能吃掉家庭开销的三成，而全国口径里教育文化娱乐的权重远没有这么重。同一个城市，同一个月，一家人喊涨，另一家人没感觉。}
  \flowitem{02}{房价不进 CPI}{CPI 量的是消费，买房在统计口径里算投资。进篮子的只有房租，外加一个叫\term{自有住房服务}{owner's equivalent rent}\index{094@自有住房服务}的虚拟项目：假设你把房子租给自己，能收多少。房价翻倍而 CPI 不动，完全可能发生。这不是造假，是定义。可对一个正在攒首付的年轻人，这定义等于把他这辈子最大的一次涨价划掉了。}
  \flowitem{03}{人对涨价敏感，对降价没记性}{她记得住排骨那 2 块 7，记不住四年前 3800 买的手机今天 2100 还更好用。早餐、打车、菜价这种高频小额支出，在记忆里的权重远高于家电和机票。行为经济学管这叫损失厌恶。人话是：疼比爽记得久。}
  \flowitem{04}{质量调整}{同样 4699 块的手机，今年性能是三年前的两倍，统计上要把这部分算成降价，术语叫\term{质量调整}{hedonic adjustment}\index{086@质量调整}。逻辑站得住。可赵敏芝并不觉得自己变富了，她只觉得手机本来就该越来越好。}
\end{flowlist}

"官方数据是编的"这个说法，中国有人讲，美国和欧洲也有人讲，讲了几十年。

我认为它是错的。CPI 的毛病不在造假，在口径，而口径是公开的、可查的、可以被批评的。值得吵的是"该不该把房价放进篮子"，不是"数字是不是编的"。

\section{三条来路}

\begin{tablewrap}
\begin{xltabular}{\linewidth}{>{\raggedright\arraybackslash}p{0.20\linewidth}XX}
\toprule
\bthead{类型} & \bthead{机制} & \bthead{典型画面} \\
\midrule
\textbf{需求拉动} & 钱多货少，总需求跑在产能前面 & 2021 年美国把支票寄进居民信箱，而工厂还没恢复 \\
\textbf{成本推动} & 产量没变，生产成本变了 & 1973 年石油危机；2026 年的关税和中东局势 \\
\textbf{预期自我实现} & 人们相信会涨，提前行动，于是真涨了 & 工人要求涨 5\%，企业提前提价 5\% \\
\bottomrule
\end{xltabular}
\end{tablewrap}

\begin{egbox}{举个栗子}
兴顺街拐角那家馄饨馆，老板姓崔。一碗三鲜馄饨 13 块，成本是：肉馅虾仁和面粉 3.9 元，一次性餐具 0.7 元，摊到每碗的房租 2.6 元，人工 3.4 元，水电燃气 1.1 元。剩 \textbf{1.3 元}归他。

猪肉和虾仁一起涨了 22\%，3.9 变成 4.76。他的利润从 1.3 掉到 \textbf{0.44 元，缩水 66\%}。

崔师傅只有两条路：提到 14 块，或者少赚三分之二。街上七八家早点铺面对同一道题，绝大多数会选前一条。

这里没有任何人多印钱，也没有任何人需求过热。物价照样涨了。

\end{egbox}

三条来路里，央行最怕第三条。前两条是外部条件变了，第三条是人的信念变了，而信念会自我实现。全社会一旦默认"每年都得涨 5\%"，这个数字就会写进工资合同、租约和供应商报价单，然后它真的成立。

央行真正在管的从来不是物价，是关于物价的信念。

信念看不见，但它有价格。

市场上同时挂着两种美国国债：一种付固定票息，一种的本金跟着 CPI 走（\term{通胀保值国债}{TIPS}\index{062@通胀保值国债}）。把两者的收益率一减，剩下那个差额，就是市场此刻押注的未来通胀，行话叫\term{盈亏平衡通胀率}{breakeven inflation}\index{077@盈亏平衡通胀率}。

央行官员每天盯的就是这条线。它一往上飘，说明"关于物价的信念"开始松动了。第 10 章会讲这类债券怎么定价。

\section{为什么不按到零}

既然通胀在偷钱，干脆按到零不行吗。

不行，零是个危险的目标。

\begin{flowlist}
  \flowitem{01}{名义工资是刚性的}{降薪极难。企业宁可裁掉 11\% 的人，也不愿给所有人普降 8\% 的工资，后者能在三天里毁掉一支队伍。而在 2\% 的温和通胀下，给某个岗位"涨 0"，他的\emph{实际}工资就自动降了 2\%，没有人需要开那个会。有点阴险，但省血。}
  \flowitem{02}{得给降息留弹药}{名义利率大致有个下限，负利率试过，副作用不小。长期通胀若是 0\%，正常时期的名义利率也会低 2 个百分点；衰退一来，能降的空间就少 2 个点。}
  \flowitem{03}{CPI 自带向上的测量偏差}{猪肉贵了赵敏芝改买鸡腿，可篮子权重更新有滞后，这叫替代偏差；新产品进篮子也慢；质量调整做不干净。统计出来的通胀通常比真实的生活成本上升高零点几个百分点。把目标定在 0，等于在瞄准负数。}
  \flowitem{04}{离通缩远一点}{目标定在 0，意味着接近一半的时间物价在跌。}
\end{flowlist}

\begin{deepbox}{进阶：可以跳过}
2\% 没有任何神圣性。

它最早被写进法律，是新西兰 1989 年的《储备银行法》。而那个数字的由来带着点偶然：财政部长先在一次电视访谈里随口说了个区间，后来才被磨成正式目标。别的国家看着好使，就抄了过去，一抄三十多年。这段掌故常被引用，细节在不同版本里略有出入。但有一点没争议：它不是从哪个方程里解出来的最优解。

学界一直在吵。一派主张提到 3\% 甚至 4\%：既然名义利率的下限卡在零附近，通胀目标高一点，正常时期的名义利率就高一点，衰退时能降的空间也大一点。另一派主张干脆换成价格水平目标制，某一年通胀低了，第二年要补回来，让物价水平本身回到轨道上，而不是只盯增速。

这件事目前没人给得出定论。

而 2026 年的现实是，主要央行都在极力避免\emph{看起来像}在放松目标。沃什那句"2\% 的通胀目标不会动摇"就是这个用途。市场一旦认为目标可以商量，前面说的盈亏平衡通胀率立刻会往上跑，届时要把它按回去的代价，比当初守住它大得多。

\end{deepbox}

第四条得展开，它最要命。核心机制是欧文·费雪 1933 年提出的\term{债务-通缩螺旋}{debt-deflation spiral}\index{083@债务-通缩螺旋}。

\begin{chainflow}
\chainnode{物价下跌}\chainarrow \chainnode{债务的实际负担上升}\chainarrow \chainnode{企业变卖资产还债}\chainarrow \chainnode{抛售压低价格}\chainarrow \chainnode{物价继续下跌}\chainarrow \chainnode{实际负担又上升}
\end{chainflow}

费雪那句总结很冷：债务人还得越多，欠得越多。他们卖资产还债的动作本身在往下砸价格，而债务的面值纹丝不动。消费者那边同时在踩刹车：反正下个月更便宜，那就等等。

\begin{egbox}{举个栗子}
赵敏芝的女婿胡建军，34 岁，在沈阳一家汽车零部件厂做设备工程师。2021 年买房贷了 94 万，现在还剩 \textbf{87.3 万}，去年到手 21.6 万。债务是他年收入的 4.04 倍。

往后十年他什么都不做，只是被动地活在两种物价环境里。

\textbf{一年 3\% 的通胀，工资同步涨。}十年后到手 21.6×1.03\textsuperscript{10} = \textbf{29.03 万}，87.3 万本金一分没涨，债务收入比降到 \textbf{3.01 倍}。

\textbf{一年 2\% 的通缩，工资也跟着降。}十年后到手 21.6×0.98\textsuperscript{10} = \textbf{17.65 万}，债务收入比升到 \textbf{4.95 倍}。

同一个人，同一笔贷款，什么都没干，负担差了 \textbf{64\%}。通缩不需要任何人犯错，就能让所有债务人同时变穷。

\end{egbox}

日本用三十年演示了这个坑有多深。泡沫 1990 年破掉之后，物价长期趴着，企业忙着还债而不是投资，居民默认明天不会更贵。日本央行 2024 年 3 月才结束负利率与 YCC，2025 年 1 月加到 0.5\%，2025 年 12 月 0.75\%，\textbf{2026 年 6 月 16 日加到 1.0\%}，是 1995 年 9 月以来的最高。

一个发达经济体花了三十年，把利率抬回 1\%。

\section{谁在偷，谁在赚}

\begin{pullquote}{}
通胀不毁灭财富，它只搬运。搬的过程没有声音，被搬的那一方通常最后才知道。
\end{pullquote}

原因在于债务合同都是用名义金额写的。胡建军欠银行 87 万 3，不管十年后一碗馄饨卖到多少钱，他还的都是 87 万 3。尺子变短，数字没变，负担自动轻了。

\begin{tablewrap}
\begin{xltabular}{\linewidth}{>{\raggedright\arraybackslash}p{0.20\linewidth}Xr}
\toprule
\bthead{谁} & \bthead{通胀来了会怎样} & \bthead{方向} \\
\midrule
\textbf{胡建军这样的房贷借款人} & 债务名义固定，收入和房价名义上涨 & 受益 \\
\textbf{赵敏芝这样把钱放定期的人} & 数字不动，购买力一年一年往下掉 & 受损 \\
\textbf{持有实物资产和股权的人} & 长期跟着名义变量走，但冲击刚发生时股债经常一起跌 & 多数受益 \\
\textbf{靠固定名义收入过活的人} & 固定票息的债券持有人、签了长约的房东、调整滞后的退休金 & 受损最重 \\
\textbf{政府} & 债务名义固定，税基跟着名义 GDP 一起膨胀 & 通常是最大赢家 \\
\bottomrule
\end{xltabular}
\end{tablewrap}

最后那一行有个名字，叫\term{没有立法的税}{taxation without legislation}\index{036@没有立法的税}。不经过任何一次投票，转移就完成了。

\begin{egbox}{举个栗子}
赵敏芝去年把铁西那套老房子卖了 96 万，全存了三年定期。她的说法是"等看清楚再说"。

按 3\% 的通胀算，十年后这笔钱能买到的东西，相当于今天的 96 ÷ 1.3439 = \textbf{71.4 万}。按 5\% 算，是 \textbf{58.9 万}。

"先放着"本身也是一个决定。它把风险从价格波动换成了购买力的慢性流失，没有把风险取消掉。

\end{egbox}

\begin{warnbox}{小心}
上面这段算术极容易被拿去当推销词。"现金会贬值，所以你必须马上买某某"是几乎所有话术的标准开场白。\hlwarm{购买力缩水是事实，"因此某项资产一定跑赢通胀"是推论，而且经常是错的推论。}本书不提供任何投资建议，第 21 章和第 23 章会专门处理。

\end{warnbox}

\section{两场相反的仗}

\begin{statrow}{3}
  \statitem{3.50 至 3.75}{\%}{美国联邦基金利率目标区间，2026 年 7 月，连续第 5 次维持不变}
  \statitem{5.02}{\%}{日本 2026 年春斗薪资涨幅预期}
  \statitem{6.3}{\%}{日本 5 月 PPI 同比}
\end{statrow}

\textbf{美国的通胀还在，而且很粘。}截至 2026 年 8 月，压力主要来自\hlmark{供给侧}：关税抬高进口成本，中东局势推高能源价格，去全球化带来重构成本，AI 资本开支拉动电力和设备需求。2025 年全球数据中心资本开支同比增长约 57\%，它们要的电和变压器都是真实且稀缺的东西。FOMC 的表述是"通胀仍然高企，部分受冲击影响"，市场上一度讨论过加息。

\begin{talkbox}
  \talkline{新闻}{美联储主席凯文·沃什 2026 年 7 月表示："通胀问题不可能在九周内解决。""2\% 的通胀目标不会动摇。"}
  \talkline{翻译}{第一句在说：这轮通胀的根子在供给侧，利率工具变不出更多的石油、芯片和电力，别指望一两次会议就有结果。第二句在说：目标短期够不着，但我绝不暗示它可以商量。一旦松口，预期那道堤坝就开始漏水，到时候要付的代价比现在大得多。}
\end{talkbox}

\textbf{日本那边，久违的螺旋居然是好消息。}2026 年春斗薪资涨幅预期约 5.02\%，5 月 PPI 同比 6.3\%。工资推物价、物价再推工资，这在别的国家是刺耳的警报，在日本几乎算一句祝福：钉了三十年的通缩心态终于松动。日本央行的任务因此换了一个，确认苗头不会变成失控。

\textbf{中国的问题是反的。}这边不是钱太多，是\textbf{需求不足、物价低迷}。货币政策基调"适度宽松"，7 天期逆回购政策利率 1.4\%，2026 年 8 月 20 日的 1 年期 LPR 3.0\%、5 年期以上 3.5\%，连续多月持平。商业银行净息差二季度末降到 1.41\% 的历史低位，这个数字本身就说明降息空间被挤得很紧。政策重心因此从"放水"挪到"让人愿意花钱"：商务部等 7 部门提出到 2030 年社零达 60 万亿元左右，8 月底又出了一套房地产组合拳。

同一个月，华盛顿在想怎么把物价按下去，东京在庆幸物价终于起来了，北京在想怎么让人愿意花钱。

教科书里那个"通胀"是单数。

\section{扔掉两个直觉}

\begin{mythbox}{常见误解}
\mvfrow{误解}{"印钱必然导致通胀。央行一放水，物价就一定涨，这是常识。"}
\mvfrow{事实}{2009 年到 2019 年，美联储、欧洲央行、日本央行做了人类历史上规模最大的\term{量化宽松}{QE}\index{033@量化宽松}，央行资产负债表扩张了好几倍，基础货币暴增。\textbf{这十年发达国家的 CPI 通胀始终低于 2\% 的目标}，欧元区和日本还长期在零附近挣扎。当年预言"恶性通胀马上来"的人排着队，等了十年也没等到。}
2021 到 2023 那一轮真正的通胀，主因是疫情造成的供应链中断和产能损失、财政把钱直接转到居民账户、需求结构从服务业剧烈切换到实物商品，以及后来的能源冲击。

关键差别在这里。

QE 不是"印钱发给老百姓"，是央行用新创造的准备金从金融机构手里买走债券。银行的资产结构变了，债券换成准备金，可你我账户里的存款没有因此多出一笔。所以 QE 推高的是资产价格，不是货架上的价格。2020 年那些寄到家里的支票是另一件事。

但别滑到另一个极端。\textbf{这不代表货币无关紧要。}

战时财政货币化、恶性通胀这类极端情形里，货币是决定性的。准确的说法是：简单的货币数量论，M 涨多少 P 就涨多少，\hlmark{已经不能用来做短期预测}。现代央行早就不拿货币总量当中介目标，改盯利率和预期。

\end{mythbox}

\begin{mythbox}{常见误解}
\mvfrow{误解}{"加了关税，进口的东西贵了，这就是通胀。"}
\mvfrow{事实}{关税制造的是\textbf{价格水平的一次性跳升}，不是通胀率的持续上升。这两件事在 2026 年格外需要分清楚。}
假设关税让某类商品一次性贵了 10\%。生效之后的 12 个月里，同比 CPI 确实会被这 10\% 抬高。可到了第 13 个月，只要关税不再加码，这个效应\textbf{自动消失}，因为做同比的基数里已经含着它了。价格停在更高的台阶上，不再往上爬。

它变成持续通胀只有一条路：\hlmark{这次涨价改变了人们的预期}，工人据此要求加薪，企业据此提前调价，第二年第三年接着走。

这就是 2026 年美联储最难受的地方：加息变不出更多的石油、芯片和电力，它能做的只有把预期按住，代价是牺牲需求和就业。沃什那句"不可能在九周内解决"不是拖延，是诚实。

\end{mythbox}

\bookbreak

回到那条主线。钱是记账，记账要有人信。通胀就是这套记账系统的计量单位在悄悄变小：它不改动任何一行数字，只改动每一行数字的意思。而债务合同按名义金额写死，单位一变，账就自己动了。

赵敏芝那 96 万还在定期里。她还在等看清楚。

\begin{takeaway}{本章一句话}
通胀不是"某样东西涨价"，是记账单位本身在缩水；它不毁灭财富，只在债权人和债务人之间做一次没人签字的再分配。

\begin{itemize}
  \item 官方通胀和体感通胀对不上，是因为篮子不同、房价不计入、人对涨价更敏感、质量调整。四条都成立，不需要阴谋论。
  \item 2\% 是安全边距不是神圣数字：给工资调整留润滑，给降息留弹药，抵消测量偏差，离通缩远一点。
  \item 通缩比通胀难治。它靠债务实际负担的上升自我强化，不需要任何人犯错。
  \item 2026 年的全球通胀主要来自供给侧，而货币政策治不了供给，只能锚住预期。
\end{itemize}

\end{takeaway}

\begin{bridgebox}
\textbf{下一章}通胀在债权人和债务人之间搬钱。可"债权"和"债务"到底是怎么记下来的？下一章给你一副眼镜，戴上它，从赵敏芝的定期到国家的财政赤字都会变得半透明。
\end{bridgebox}

